% =====================================================================
%  Section 1 -- Introduction (owner: LEAD; revised by TEXT-REVISER).
%  Figure 1 (fig:design) and Table 1 (tab:notation) are \input here so that
%  they float next to their first mention. Facts in the first paragraph are
%  documented in bib/lead.bib.
% =====================================================================
\section{Introduction}\label{sec:intro}

Recommender systems decide much of what the users of large online services
see. The European Union's Digital Services Act (DSA) puts it in its recitals:
such systems ``play an important role in the amplification of certain
messages, the viral dissemination of information and the stimulation of
online behaviour'' \citep[Recital~70]{eu2022dsa}. For the largest services,
those with at least 45~million average monthly active recipients in the
Union, the Act requires providers to assess the systemic risks that stem from
their services, among them negative effects on human dignity and
non-discrimination, taking the design of their recommender systems into
account (Articles~33 and~34). For Jewish users the question is not abstract.
In the third survey of Jewish people by the EU Agency for Fundamental Rights,
conducted online in 2023, 90\% of almost 8{,}000 respondents in 13 Member
States said they had encountered antisemitism online in the previous twelve
months \citep{fra2024antisemitism}. That figure says that antisemitic content
is widely seen. It does not say why.

There are two broad answers, and they call for different remedies. Users may
see antisemitic content because they, or the accounts they follow, seek it
out: the exposure is driven by \emph{demand}. Or a recommender may put such
content in front of them more often than an alternative way of choosing
content would: the exposure is driven by \emph{supply}, and the remedy lies in
the design of the recommender. Public argument about whether algorithms
``amplify'' hateful content rarely says what amplification is compared with,
or for whom, and published audits differ in what they compare, what they
measure, the content they study and the period they cover
(\cref{sec:background}). A claim that an algorithm amplifies antisemitic
content, or that it does not, often cannot be checked against another,
because the two claims concern different estimands.

This paper supplies the estimands and the methods to estimate them.
\emph{Amplification} is a causal contrast: the expected exposure to
antisemitic content under the algorithmic policy, divided by the expected
exposure under a named reference policy, for a fixed demand profile. An audit
with sock puppets (synthetic accounts whose behaviour the researcher scripts)
then becomes a randomised experiment whose identifying assumptions can be
stated and, in part, checked. Its design problems are those of any experiment
on a live system that learns: puppets can change the model they audit, the
model changes while it is audited, the outcome is measured by an imperfect
classifier, and every session adds to the footprint of the audit on a real
service. \Cref{fig:design} shows the audit design that results.

\IfFileExists{figures/fig-design.tex}{% =====================================================================
%  Figure: the audit design, from personas to a stopping time (LEAD).
%  Glyphs are SCHEMATIC: every number drawn (cell counts, slates, score
%  densities, bar heights, intervals, the confidence sequence) is invented
%  for legibility and is neither a simulation output nor data.
%  Requires the house colours defined in preamble.tex.
% =====================================================================
\begin{figure}[t]
\centering
\begin{tikzpicture}[x=1cm,y=1cm,
    tile/.style={rounded corners=3pt,draw=LabRule,line width=0.6pt,fill=white,
                 minimum width=3.5cm,minimum height=3.0cm,inner sep=0pt},
    ttl/.style={font=\footnotesize\bfseries,text=LabNavy,align=center,
                anchor=north,inner sep=0pt,text width=3.35cm},
    ref/.style={font=\scriptsize,text=LabGrey,align=center,anchor=south,
                inner sep=0pt,text width=3.35cm},
    flowarrow/.style={-{Stealth[length=4.5pt]},LabNavy,line width=0.8pt},
    lab/.style={font=\tiny,text=LabGrey,inner sep=0.5pt},
    algc/.style={fill=LabNavy},
    refc/.style={fill=LabNavyLight},
  ]
  % ------------------------------------------------------------------
  % tile helper:  \tile{name}{x}{y}{title}{reference}
  \newcommand{\tile}[5]{%
    \node[tile] (#1) at (#2,#3) {};
    \node[ttl] at ([yshift=-5pt]#1.north) {#4};
    \node[ref] at ([yshift=4pt]#1.south) {#5};}
  \def\yA{0}\def\yB{-4.0}
  \def\gx{-0.27}  % vertical offset of glyph centres within a tile

  % ---------------- row I: from design to labelled exposures -----------
  \tile{t1}{1.75}{\yA}{Personas $\times$\\policies}%
    {\S\ref{sec:prot-personas}, Def.~\ref{def:amplification}}
  \tile{t2}{5.8}{\yA}{Randomised puppets,\\blocked by wave}%
    {Thm.~\ref{thm:identify}, Prop.~\ref{prop:interference}}
  \tile{t3}{9.85}{\yA}{Training, then\\measurement sessions}%
    {\S\ref{sec:prot-personas}, Prop.~\ref{prop:markov}}
  \tile{t4}{13.9}{\yA}{Frozen classifier\\labels every item}%
    {\S\ref{sec:prot-validation}}

  % ---------------- row II: from labels to a decision ------------------
  \tile{t5}{1.75}{\yB}{Validation sample,\\coded blind}%
    {\S\ref{sec:prot-validation}}
  \tile{t6}{5.8}{\yB}{Corrected exposure\\in each arm}%
    {Thm.~\ref{thm:misclass}}
  \tile{t7}{9.85}{\yB}{Amplification\\for each persona}%
    {Def.~\ref{def:amplification}, Thm.~\ref{thm:decomp}}
  \tile{t8}{13.9}{\yB}{Confidence sequence\\and stopping time}%
    {Alg.~\ref{alg:sequential}, Thm.~\ref{thm:sequential}, Prop.~\ref{prop:power}}

  % ---------------- arrows ----------------
  \foreach \a/\b in {t1/t2,t2/t3,t3/t4,t5/t6,t6/t7,t7/t8}{
    \draw[flowarrow] (\a.east) -- (\b.west);}
  \draw[flowarrow,rounded corners=6pt] (t4.south) -- ++(0,-0.5) -| (t5.north);
  \node[fill=white,font=\scriptsize,text=LabNavy,inner sep=3pt] at (7.825,-2.0)
    {after every wave: new items are labelled, a validation subsample is coded, and
     every estimate is updated};
  % the sequential loop: back to the next wave unless a boundary is crossed
  \draw[flowarrow,LabAccent,rounded corners=6pt,dash pattern=on 3pt off 1.6pt]
    (t8.east) -- ++(0.32,0) |- ([yshift=0.42cm]t2.north) -- (t2.north);
  \node[fill=white,font=\scriptsize,text=LabAccent,inner sep=2pt,anchor=west]
    at (6.3,1.92) {next wave, unless the sequence has crossed a decision boundary
    (\S\ref{sec:prot-stopping})};

  % ================== glyphs ==================
  % 1. personas x policies: |D| x 2 cells, puppets in each
  \begin{scope}[shift={(t1.center)},yshift=\gx cm]
    \def\cw{0.5}\def\ch{0.25}
    \foreach \c/\st/\nm in {0/algc/alg,1/refc/ref}{
      \node[lab,anchor=south] at ({-0.42+\c*\cw+\cw/2},0.52) {\nm};
      \foreach \r in {0,1,2,3}{
        \fill[LabTint] ({-0.42+\c*\cw},{0.5-\r*\ch}) rectangle ++(\cw,-\ch);
        \foreach \k in {0,1,2}{
          \fill[\st] ({-0.42+\c*\cw+0.11+\k*0.14},{0.5-\r*\ch-\ch/2}) circle (1.45pt);}
      }}
    \draw[LabRule,line width=0.3pt,step=\cw] (-0.42,-0.5) grid (0.58,0.5);
    \draw[LabRule,line width=0.3pt] (-0.42,0.25) -- (0.58,0.25) (-0.42,0) -- (0.58,0)
      (-0.42,-0.25) -- (0.58,-0.25);
    \draw[LabNavy,line width=0.5pt] (-0.42,-0.5) rectangle (0.58,0.5);
    \foreach \r in {0,1,2,3}{
      \node[lab,anchor=east] at (-0.44,{0.5-\r*\ch-\ch/2}) {$d_{\r}$};}
    \node[lab,anchor=west,align=left] at (0.66,0.0) {persona\\[-1pt]$d\in\mathcal D$\\[1pt]policy\\[-1pt]$\pi$};
  \end{scope}

  % 2. randomised puppets in waves (equal numbers per arm; one sentinel each)
  \begin{scope}[shift={(t2.center)},yshift=\gx cm+0.08cm]
    \foreach \w/\xo/\pat in {1/-1.38/{1,0,0,1,0,1,1,0},2/-0.74/{0,1,1,0,1,0,0,1},
                             3/-0.10/{1,1,0,0,0,1,0,1}}{
      \draw[LabRule,fill=LabTint,rounded corners=2pt,line width=0.4pt]
        (\xo,-0.5) rectangle ++(0.56,1.0);
      \foreach \a [count=\j from 0] in \pat {
        \pgfmathtruncatemacro{\cx}{mod(\j,2)}\pgfmathtruncatemacro{\cy}{div(\j,2)}
        \ifnum\a=1 \def\fc{LabNavy}\else\def\fc{LabNavyLight}\fi
        \fill[\fc] ({\xo+0.16+0.24*\cx},{0.36-0.2*\cy}) circle (1.6pt);}
      \draw[LabNavy,line width=0.5pt,fill=white] ({\xo+0.28},-0.4) circle (1.6pt);
      \node[lab,anchor=north] at ({\xo+0.28},-0.52) {$w{=}\w$};}
    \node[font=\scriptsize,text=LabGrey] at (0.66,0.02) {$\cdots$};
    \node[lab,anchor=west,align=left] at (0.86,0.36) {\textcolor{LabNavy}{$\bullet$}\,alg};
    \node[lab,anchor=west,align=left] at (0.86,0.16) {\textcolor{LabNavyLight}{$\bullet$}\,ref};
    \node[lab,anchor=west,align=left] at (0.86,-0.4) {$\circ$\,sentinel};
  \end{scope}

  % 3. training, then measurement: slates of k positions over T steps
  \begin{scope}[shift={(t3.center)},yshift=\gx cm+0.12cm]
    % training: seed items
    \foreach \i in {0,...,3}{
      \fill[LabNavyLight!45] (-1.42,{0.28-0.2*\i}) rectangle ++(0.26,-0.14);}
    \node[lab,anchor=south] at (-1.29,0.36) {train};
    \draw[-{Stealth[length=2.5pt]},LabGrey,line width=0.5pt] (-1.06,0) -- (-0.8,0);
    % measurement: columns are steps t, rows slate positions; one item opened per step
    \def\cw{0.27}
    \foreach \col [count=\t from 0] in {{0,1,0},{0,0,2},{1,0,0},{2,0,1},{1,1,0},{0,2,0},{1,0,0}}{
      \foreach \s [count=\r from 0] in \col {
        \ifcase\s \def\fc{LabNavyLight!45}\or\def\fc{LabNavyMid}\or\def\fc{LabAccent}\fi
        \fill[\fc] ({-0.7+\t*\cw},{0.3-\r*0.24}) rectangle ++(0.21,-0.18);}}
    % opened item per step (never a flagged one)
    \foreach \t/\r in {0/0,1/0,2/1,3/1,4/0,5/0,6/1}{
      \draw[LabNavy,line width=0.7pt] ({-0.7+\t*\cw-0.03},{0.3-\r*0.24+0.03})
        rectangle ++(0.27,-0.24);}
    \node[lab,anchor=south] at (0.17,0.36) {measure: slate $\times$ step $t$};
    \node[lab,anchor=north,align=center] at (0.17,-0.42)
      {\textcolor{LabNavy}{$\square$}\,opened\enspace
       \textcolor{LabAccent}{$\blacksquare$}\,flagged:\\[-1pt]recorded, never opened};
  \end{scope}

  % 4. frozen classifier: score densities of L = 0 and L = 1, one threshold
  \begin{scope}[shift={(t4.center)},yshift=\gx cm]
    \fill[LabNavyLight!40,domain=0.05:0.9,samples=30,smooth]
      plot (\x,{-0.42+0.85*exp(-((\x+0.45)^2)/0.13)}) -- (0.9,-0.42) -- (0.05,-0.42) -- cycle;
    \fill[LabAccent!25,domain=-1.15:0.05,samples=30,smooth]
      plot (\x,{-0.42+0.6*exp(-((\x-0.62)^2)/0.16)}) -- (0.05,-0.42) -- (-1.15,-0.42) -- cycle;
    \draw[LabNavyLight,line width=0.8pt,domain=-1.25:0.9,samples=50,smooth]
      plot (\x,{-0.42+0.85*exp(-((\x+0.45)^2)/0.13)});
    \draw[LabAccent,line width=0.8pt,domain=-0.6:1.3,samples=50,smooth]
      plot (\x,{-0.42+0.6*exp(-((\x-0.62)^2)/0.16)});
    \draw[LabNavy,line width=0.4pt] (-1.3,-0.42) -- (1.35,-0.42);
    \draw[LabNavy,line width=0.6pt,dash pattern=on 2pt off 1pt] (0.05,-0.48) -- (0.05,0.5);
    \node[lab,text=LabNavyLight,anchor=south] at (-0.45,0.44) {$L_c{=}0$};
    \node[lab,text=LabAccent,anchor=south] at (0.72,0.2) {$L_c{=}1$};
    \node[lab,anchor=north] at (0,-0.44) {score; $\widehat L_c=1$ right of the threshold};
  \end{scope}

  % 5. validation sample: score bands, high bands oversampled; blind double coding
  \begin{scope}[shift={(t5.center)},yshift=\gx cm]
    \foreach \b/\n in {0/2,1/2,2/3,3/5,4/6}{
      \pgfmathsetmacro{\sh}{8+9*\b}
      \fill[LabNavy!\sh] (-1.05,{-0.5+0.2*\b}) rectangle ++(1.5,0.18);
      \foreach \k in {1,...,\n}{
        \pgfmathsetmacro{\xx}{-1.05+1.5*\k/(\n+1)}
        \pgfmathtruncatemacro{\isa}{(\b>2 && mod(\k,3)!=0) || (\b==2 && \k==2) ? 1 : 0}
        \ifnum\isa=1 \def\fc{LabAccent}\else\def\fc{LabNavy}\fi
        \fill[\fc] (\xx,{-0.41+0.2*\b}) circle (1.3pt);}}
    \node[lab,anchor=east] at (-1.07,0.4) {high};
    \node[lab,anchor=east] at (-1.07,-0.4) {low};
    % two coders and an adjudicator
    \foreach \yy in {0.3,-0.05}{
      \fill[LabNavyMid] (0.85,\yy+0.07) circle (1.6pt);
      \fill[LabNavyMid] (0.85,\yy-0.07) ellipse (2.6pt and 1.8pt);}
    \fill[LabNavy] (1.15,0.2) circle (1.6pt);
    \fill[LabNavy] (1.15,0.06) ellipse (2.6pt and 1.8pt);
    \node[lab,anchor=north,align=center] at (1.0,-0.2) {2 coders\\[-1pt]+ adjudicator};
  \end{scope}

  % 6. corrected exposure in each arm: naive (outlined) versus corrected (solid)
  \begin{scope}[shift={(t6.center)},yshift=\gx cm]
    \draw[LabNavy,line width=0.4pt] (-1.2,-0.45) -- (1.25,-0.45);
    \foreach \x/\nv/\cv/\c/\nm in {-0.7/0.78/0.62/LabNavy/alg,0.45/0.40/0.20/LabNavyLight/ref}{
      \draw[\c,line width=0.6pt,fill=white] (\x,-0.45) rectangle ++(0.3,\nv);
      \fill[\c] (\x+0.36,-0.45) rectangle ++(0.3,\cv);
      \node[lab,anchor=north] at (\x+0.33,-0.47) {\nm};}
    \draw[LabNavyMid,line width=0.5pt,fill=white] (0.12,0.4) rectangle ++(0.13,0.1);
    \node[lab,anchor=west] at (0.27,0.45) {naive};
    \fill[LabNavyMid] (0.12,0.22) rectangle ++(0.13,0.1);
    \node[lab,anchor=west] at (0.27,0.27) {corrected};
  \end{scope}

  % 7. amplification for each persona: log AR with interval, margin [1/delta, delta]
  \begin{scope}[shift={(t7.center)},yshift=\gx cm]
    \def\lu{1.15}  % cm per unit of log AR
    \fill[LabAccent!12] ({-0.223*\lu+0.1},-0.5) rectangle ({0.223*\lu+0.1},0.5);
    \draw[LabAccent,line width=0.4pt,dash pattern=on 1.5pt off 1pt]
      ({-0.223*\lu+0.1},-0.5) -- ({-0.223*\lu+0.1},0.5)
      ({0.223*\lu+0.1},-0.5) -- ({0.223*\lu+0.1},0.5);
    \draw[LabNavy,line width=0.5pt] (0.1,-0.5) -- (0.1,0.5);
    \foreach \r/\est/\lo/\hi in {0/0.05/-0.15/0.18,1/0.40/0.25/0.58,2/0.20/-0.05/0.47,3/0.85/0.48/1.15}{
      \pgfmathsetmacro{\yy}{0.38-0.25*\r}
      \draw[LabNavyLight,line width=1.6pt] ({\lo*\lu+0.1-0.06},\yy) -- ({\hi*\lu+0.1+0.06},\yy);
      \draw[LabNavy,line width=1.0pt] ({\lo*\lu+0.1},\yy) -- ({\hi*\lu+0.1},\yy);
      \fill[LabNavy] ({\est*\lu+0.1},\yy) circle (1.7pt);
      \node[lab,anchor=east] at (-1.25,\yy) {$d_{\r}$};}
    \node[lab,anchor=north] at (0.1,-0.52) {$\log\mathrm{AR}(d)$;\ \textcolor{LabAccent}{margin $\pm\log\delta$}};
  \end{scope}

  % 8. confidence sequence for log AR(d) over waves; stops at tau
  \begin{scope}[shift={(t8.center)},yshift=\gx cm]
    \def\wx{0.29}   % cm per wave
    \def\vy{0.72}   % cm per unit of log AR
    \def\yo{-0.4}   % y of log AR = 0
    \draw[LabNavy,line width=0.4pt,-{Stealth[length=2.5pt]}] (-1.15,\yo) -- (1.35,\yo);
    \draw[LabNavy,line width=0.4pt] (-1.15,-0.52) -- (-1.15,0.55);
    % margin log(delta) = 0.223
    \draw[LabAccent,line width=0.5pt,dash pattern=on 2pt off 1pt]
      (-1.15,{\yo+0.223*\vy}) -- (1.3,{\yo+0.223*\vy});
    % the band (waves 1..4), then the stop
    \fill[LabNavy!18]
      plot coordinates {(-1.15+1*\wx,{\yo+1.40*\vy}) (-1.15+2*\wx,{\yo+0.95*\vy})
                        (-1.15+3*\wx,{\yo+0.86*\vy}) (-1.15+4*\wx,{\yo+0.88*\vy})}
      -- plot coordinates {(-1.15+4*\wx,{\yo+0.27*\vy}) (-1.15+3*\wx,{\yo+0.17*\vy})
                        (-1.15+2*\wx,{\yo+0.03*\vy}) (-1.15+1*\wx,{\yo-0.20*\vy})} -- cycle;
    \draw[LabNavy,line width=0.6pt] plot coordinates {(-1.15+1*\wx,{\yo+0.60*\vy})
      (-1.15+2*\wx,{\yo+0.49*\vy}) (-1.15+3*\wx,{\yo+0.52*\vy}) (-1.15+4*\wx,{\yo+0.58*\vy})};
    % waves not run
    \foreach \w in {5,...,8}{
      \draw[LabRule,line width=0.4pt] (-1.15+\w*\wx,\yo) -- ++(0,-0.05);}
    \foreach \w in {1,...,4}{
      \draw[LabNavy,line width=0.4pt] (-1.15+\w*\wx,\yo) -- ++(0,-0.05);}
    \draw[LabAccent,line width=0.7pt] (-1.15+4*\wx,\yo-0.07) -- (-1.15+4*\wx,{\yo+0.27*\vy});
    \fill[LabAccent] (-1.15+4*\wx,{\yo+0.27*\vy}) circle (1.3pt);
    \node[lab,text=LabAccent,anchor=north] at (-1.15+4*\wx,\yo-0.06) {$\tau$};
    \node[lab,text=LabAccent,anchor=west] at (0.1,{\yo+0.223*\vy+0.1}) {$\log\delta$};
    \node[lab,anchor=north] at (0.62,\yo-0.06) {waves};
    \node[lab,anchor=east] at (-1.17,{\yo}) {$0$};
  \end{scope}
\end{tikzpicture}
\caption{\textbf{The audit design, from personas to a stopping time.} Top row:
each persona $d$ (a fixed demand profile) is crossed with the algorithmic and
a reference policy; puppets are assigned to the $|\mathcal D|\times2$ cells
completely at random within waves, which are blocks, and untrained sentinels
monitor interference and drift. Every puppet is trained in the same way; the
policy contrast starts with the measurement sessions, in which each slate is
recorded, every item is scored by a classifier frozen at registration, and the
puppet opens one item per step but never a flagged one. Bottom row: a
stratified, oversampled validation sample is coded blind to arm, persona and
score, and gives the sensitivity and specificity of the classifier in each
arm; exposure is corrected arm by arm, because error rates can differ when the
content mix differs; the amplification ratio is estimated for each persona and
compared with a registered margin $[1/\delta,\delta]$; and a confidence
sequence for $\mathrm{AR}(d)$, built from each puppet's confirmed flagged
exposure and (asymptotically) valid whenever the audit stops, decides after
each wave whether to stop. Here, on the log scale, it stops at wave $\tau=4$,
when it lies entirely above $\log\delta$ (rule S1). Labels under each tile point to the corresponding results.
\textbf{Schematic:} all glyph values are invented for legibility.}
\label{fig:design}
\end{figure}
}{}%
\paragraph{What this paper does, and does not, do}
This is a methods paper. We define the estimands, prove what a sock-puppet
audit can and cannot identify, derive estimators that correct for
misclassification and a sequential procedure that stops when the evidence is
clear, evaluate them in a simulation study on a synthetic recommender with
known ground truth, and pre-register a protocol for an audit, with its ethical
and legal conditions. We have run \emph{no} audit of any real platform, and
nothing here should be read as a finding that any service amplifies
antisemitic content, or that it does not. Published audits are cited only as
their authors report them. The protocol is platform-neutral: it is written for
any service with a recommendation surface and a reference policy. Notation is
collected in \cref{tab:notation}.

\IfFileExists{tables/tab-notation.tex}{% =====================================================================
%  Table: notation (LEAD). Symbols follow SPEC.md and Sections 3--6.
% =====================================================================
\begin{table}[t]
\caption{\textbf{Notation.} Symbols used throughout. Vectors are columns;
$\bm1$ is the vector of ones and $\bm1_{\mathcal A}$ the indicator of the
antisemitic states.}
\label{tab:notation}
\footnotesize
\renewcommand{\arraystretch}{1.1}
\newcommand{\nhead}[1]{\multicolumn{2}{@{}l}{\footnotesize\scshape\bfseries\color{LabNavy}#1}\\[1pt]}
\begin{minipage}[t]{0.49\textwidth}
\begin{tabularx}{\linewidth}[t]{@{}>{$}l<{$}@{\hspace{0.7em}}>{\raggedright\arraybackslash}X@{}}
\toprule
\nhead{Accounts, demand and policies}
u=1,\dots,U & puppet (synthetic account) \\
d\in\mathcal D,\ d_0 & persona (fixed demand profile); baseline persona \\
\mathrm{alg};\ \rho & algorithmic policy; reference ($\mathrm{chron}$, $\mathrm{rand}$, $\mathrm{pop}$) \\
\mathcal C_t & eligible pool of items at time $t$ \\
t=1,\dots,T;\ k & step of a session; slate positions recorded \\
\midrule
\nhead{Design and inference}
w=1,\dots,W;\ n_{d\pi} & wave (block); puppets in cell $(d,\pi)$ \\
\ell_w,\ \bar\ell & load: puppets' share of the engagement signal; its average \\
\xi_\pi,\ \xi_\Delta & change in exposure, and in $\mathrm{AA}$, per unit of load (\cref{prop:interference}) \\
C_t,\ [L,U] & confidence sequence after $t$ puppet pairs; its running intersection (\cref{alg:sequential}) \\
\delta;\ \alpha & registered margin; error level \\
\midrule
\nhead{Measurement}
L_c,\ \widehat L_c & true label of item $c$ ($1$: antisemitic); classifier label \\
\mathrm{Se}_\pi(d),\ \mathrm{Sp}_\pi(d) & sensitivity and specificity on the items shown in cell $(d,\pi)$ \\
J;\ \lambda_S & Youden's index $\mathrm{Se}+\mathrm{Sp}-1$; ratio $\mathrm{Se}_{\mathrm{alg}}/\mathrm{Se}_{\mathrm{ref}}$ \\
q_\pi(d) & flagged share: recorded items with $\widehat L_c=1$ \\
n_0,\ n_1 & coded validation items with $L_c=0$, with $L_c=1$ \\
\bottomrule
\end{tabularx}
\end{minipage}\hfill
\begin{minipage}[t]{0.49\textwidth}
\begin{tabularx}{\linewidth}[t]{@{}>{$}l<{$}@{\hspace{0.7em}}>{\raggedright\arraybackslash}X@{}}
\toprule
\nhead{Exposure and amplification}
E_u(\pi) & share of antisemitic items in a session under policy $\pi$ \\
\mu_\pi(d),\ \mu_\pi(Q) & mean exposure at persona $d$; averaged over $Q$ \\
\mathrm{AR}_\rho,\ \mathrm{AA}_\rho & amplification ratio; additive amplification (\cref{def:amplification}) \\
F_b & real users' demand, formed under policy $b$ \\
\tau;\ \tau_{\mathrm{sup}},\ \tau_{\mathrm{dem}},\ \iota & total effect; supply, demand-shaping and interaction terms (\cref{thm:decomp}) \\
\phi_{\mathrm{sup}},\ \phi_{\mathrm{dem}} & Shapley shares of supply and demand \\
\midrule
\nhead{Markov model of exposure}
\mathcal X\supseteq\mathcal A;\ N,B,A & content states; antisemitic states; neutral, borderline, antisemitic \\
P_{\mathrm{alg}},\ P_{\mathrm{ref}};\ \alpha_\pi & transition matrices; initial distribution \\
\varpi;\ \varpi(\mathcal A) & stationary distribution; long-run exposure \\
e_T(P,\alpha) & expected exposure in a $T$-item session \\
Z;\ H=Z-\bm1\varpi^{\top} & fundamental matrix $(I-P+\bm1\varpi^{\top})^{-1}$; group inverse of $I-P$ \\
\bm g_{\mathcal A}=H\bm1_{\mathcal A};\ \kappa_{\mathcal A} & exposure potential; sensitivity coefficient $\tfrac12\operatorname{osc}(\bm g_{\mathcal A})$ \\
M_{xy};\ K(P) & mean first-passage time; Kemeny's constant \\
\bottomrule
\end{tabularx}
\end{minipage}
\end{table}
}{}%
\paragraph{Contributions}
Classical results are attributed where used; we flag what is, to our
knowledge, new in the form given.
\begin{enumerate}[label=\textbf{C\arabic*},ref=C\arabic*,leftmargin=2.4em,itemsep=2pt]
  \item\label{contrib:estimands} \textbf{Estimands.} The amplification ratio
  $\mathrm{AR}_\rho(d)$ and additive amplification $\mathrm{AA}_\rho(d)$
  compare mean exposure under the algorithm and under a reference policy
  $\rho$ at fixed demand $d$ (\cref{def:amplification};
  \citealp{rubin1974,imbensrubin2015}); the reference is part of the
  estimand. For real users, observed exposure splits into demand, supply,
  demand shaping and a feedback interaction, bounded by the heterogeneity of
  supply times the shift in demand (\cref{thm:decomp}), and fixed-demand
  audits identify only supply (\cref{prop:audit-scope}). The bilinear form,
  the interaction bounds and \cref{prop:audit-scope} are new in this form.
  \item\label{contrib:markov} \textbf{A Markov model of exposure.} Long-run
  and finite-horizon exposure, hitting times and \emph{rabbit holes} follow
  from the fundamental matrix (\cref{prop:markov};
  \citealp{kemeny1960,meyer1975}). Redirecting a share $\varepsilon$ of
  recommendations to an antisemitic state raises long-run exposure, to first
  order, by the relative amount $\varepsilon$ times the mean number of items
  needed to reach it from equilibrium, while within a $T$-item session the
  effect is capped at $(T-1)/4$ times the change (\cref{thm:perturb}); these
  bounds are new in this form.
  \item\label{contrib:identify} \textbf{Identification.} Randomising puppets
  to persona--policy cells within waves identifies $\mathrm{AR}(d)$,
  $\mathrm{AA}(d)$ and a persona-contrast decomposition under four stated
  conditions, with the puppet as the unit of variance (\cref{thm:identify}).
  Treating waves as groups of partial interference \citep{hudgens2008} and the
  puppets' share of engagement as an exposure mapping \citep{aronow2017}, we
  bound the bias by a constant $\xi$ times that share, show that the direction
  of the bias depends on how much puppet engagement reaches false negatives,
  so no sign can be assumed, and give a randomised-saturation design and a
  breakdown value for sensitivity analysis (\cref{prop:interference}).
  \item\label{contrib:misclass} \textbf{Misclassification-corrected
  measurement.} The naive ratio of flagged shares is attenuated towards one
  through the classifier's specificity, and when error rates depend on content
  type and the arms show different mixtures of types, the naive and even a
  pooled-corrected ratio can lie on the wrong side of one
  (\cref{thm:misclass}). Validation by arm, as double sampling within puppets
  provides, is consistent in either case \citep{rogangladen1978,neyman1938}.
  \item\label{contrib:sequential} \textbf{Sequential auditing and its cost.} A
  confidence sequence from a puppet-level outcome is approximately valid at
  any data-dependent stopping time, and a betting variant exactly
  (\cref{alg:sequential,thm:sequential};
  \citealp{wald1945,howard2021,ramdas2023}). A closed-form power calculation
  (\cref{prop:power}) shows the cost: under plausible inputs a materiality
  audit of four personas needs about $4{,}900$ accounts and $170{,}000$ coded
  exposures.
  \item\label{contrib:sim} \textbf{A simulation study} on a synthetic
  recommender with known amplification evaluates these tools
  (\cref{sec:simulation}). Simulating the registered analysis itself, we find
  its normal-mixture sequences conservative, its closed-form variance
  $1.4$--$2.0$ times too large, and its exact fallback nearly powerless;
  slate-state Markov chains need a correction identifiable only for the
  two-state chain, so Markov summaries are exploratory.
  \item\label{contrib:protocol} \textbf{A pre-registered protocol with a
  feasibility gate.} We fix in advance the access route, personas,
  randomisation, interference controls, validation, estimands and decision
  rules (\cref{sec:protocol}), with ethical safeguards (\cref{sec:ethics}):
  puppets never post, interact or open flagged items. Because the materiality
  audit is infeasible for a small lab, the protocol commits only to a pilot
  and a registered feasibility gate, which under illustrative caps selects
  the detection of amplification for a single persona.
\end{enumerate}

\paragraph{What audits can and cannot show}
An audit with puppets of fixed demand identifies supply effects at the
personas it audits. It says nothing about how the algorithm shapes demand:
without outside information on how demand forms, the identified set for the
algorithm's total effect on real users is as wide as the range of exposures
under the reference policy (\cref{prop:audit-scope}). An audit can establish
that, for a stated service, surface, period and persona, the algorithm
supplies more, or less, antisemitic content than a stated alternative would
supply to the same demand. It cannot establish that the algorithm made anyone
want such content. Stating this limit is part of the method, not a disclaimer
added to it.

\paragraph{Related work and the series}
\Cref{sec:background} reviews published audits of recommender systems
\citep{sandvig2014,ribeiro2020,huszar2022,haroon2023,guess2023,hosseinmardi2024},
the supply-and-demand debate \citep{munger2022}, feedback loops
\citep{chaney2018,jiang2019}, the regulatory setting, and classifier-based
measurement. Our estimands build on mediation analysis
\citep{robins1992,pearl2001,vanderweele2014} and on the decomposition of
\citet{kitagawa1955}. This is the fourth paper in a series. The first
estimates how many antisemitic hate crimes go unrecorded
\citep{kornsteinlab2026darkfigure}, the second models surges in incidents as
self-exciting processes \citep{kornsteinlab2026contagion}, and the third maps
antisemitic attitudes \citep{kornsteinlab2026attitudes}. The first three study
what is recorded and what is believed; this one studies what is shown.

\paragraph{Roadmap}
\Cref{sec:estimands} defines the estimands and the decomposition,
\cref{sec:markov} develops the Markov model, \cref{sec:identification} gives
identification under interference and drift, and \cref{sec:measurement} the
misclassification correction, the sequential procedure and the power
calculation. \Cref{sec:simulation} reports the simulation study,
\cref{sec:protocol} summarises the registered protocol, \cref{sec:ethics}
sets out the ethical and legal conditions, and \cref{sec:discussion}
discusses limitations and extensions. Proofs are in
\cref{app:proofs-estimands,app:proofs-design}, simulation details in
\cref{app:simdetails}, and the full registered protocol in \cref{app:protocol}.
